\chapter{记忆消除的计算边界：可达性、可逆性与能力保持}
\label{chap:geometry_of_oblivion1}
上一章区分了解绑定、降权、压缩、再估计、剪枝与授权删除，本章进一步追问：当治理者要求“消除一段记忆”时，目标究竟位于当前工作态、事件记录、结构参数、外部副本，还是它们的派生关系中？若只让普通查询不再返回目标，系统可能仍能被提示恢复；若直接改写共享参数，又可能同时破坏无关能力。我们因此把软擦除、硬擦除、复发、成瘾与安全对齐统一为范围声明、依赖追踪、干预强度和能力保持问题，分别讨论可访问性改变、内容删除、参数遗忘及行为约束的证据门槛。物理图景保留为辅助直觉，但不存在由“曲率归零”自动推出安全遗忘的结论。

\section{消除对象的本体：记录、结构、索引与派生副本}
\label{sec:ontology_of_memory}
“删除记忆”若没有对象清单，就不是可执行命令。原章把短时记忆写成瞬时场与全息缓冲，把长时记忆写成度量张量的永久形变，再以负熵解释智能积累。这种二分把活动、事件、技能、索引和实际存储混成了形与质，无法说明删除请求影响哪一层。我们定义系统在离散时刻 $k$ 的记忆状态为
\[
\mathfrak M_k=(H_k,E_k,\Theta_k,I_k,X_k,D_k),
\]
其中 $H_k$ 是当前工作态与临时缓存，$E_k$ 是带身份、时间和来源的事件记录，$\Theta_k$ 是预测、策略与技能的慢结构，$I_k$ 是检索索引和路由规则，$X_k$ 是外部记忆及备份，$D_k$ 是从上述对象生成的摘要、嵌入、缓存、报告和训练派生物。各分量具有不同存续时间、权限和删除接口，不能用统一振幅或曲率表达。

每个对象写成带谱系的记录
\[
z_i=(\iota_i,\mathrm{type}_i,\mathrm{content}_i,\mathrm{source}_i,
\mathrm{owner}_i,\mathrm{policy}_i,\mathrm{version}_i,\mathrm{hash}_i),
\]
并以有向依赖图 $\mathcal G_k=(V_k,A_k)$ 记录“由何对象生成”。若摘要 $z_b$ 来自事件 $z_a$，索引 $z_c$ 又由摘要生成，则删除 $z_a$ 的请求必须判断 $z_b,z_c$ 是否也在范围内；否则原件虽消失，摘要和向量索引仍可能泄露相关信息。相反，某些派生规则由大量样本共同支持，不能仅因一个事件被删除就无条件移除整条规则。系统需要估计单个来源的可辨识影响，并把不可确定部分写进删除声明。

我们把消除请求表示为
\[
\rho=(Q,\Omega,\ell,t_d,\mathcal R,\mathcal P),
\]
其中 $Q$ 是目标身份谓词，$\Omega$ 是允许操作的存储域，$\ell$ 是所要求的消除等级，$t_d$ 是期限，$\mathcal R$ 是保留义务，$\mathcal P$ 是授权证明。消除等级至少包括：在线清空、默认不可访问、逻辑删除、物理删除、派生清理、参数影响移除和行为禁止。它们不是从软到硬的单一尺度。例如行为禁止可以很强，却并未删除知识；物理删除外部文件也可能不改变已经学习的参数；清空$H_k$只结束当前激活，不等于事件或技能消失。

工作态的历史窗口也需要准确命名。$H_k$可以包含当前候选、工具返回、尚未提交的动作和短期缓存；它们由容量、任务边界或超时规则释放。$E_k$保存的并非“全息录像”，而是经过观测、压缩和来源标记的事件，因此本来就不保证复现全部过去。$\Theta_k$中的技能表现可能由多个事件共同形成，也可能通过预训练、程序或外部工具获得。看到熟悉线索后重新产生相似反应，最多说明当前输入与残余结构共同生成了读出，不能证明存在测地线必然复活旧波形。

结构与内容的联合仍然重要，但必须通过身份接口完成。全局分布表示可以发现与目标相似的候选，局部结构网络负责核验人物、时间、来源、权限和派生关系。仅按向量相似度扩大删除范围会误伤同名人物、近义概念和合法科学知识；仅按精确键删除又会漏掉改写摘要和跨模态副本。我们采用两阶段流程：先由分布检索形成候选闭包，再由结构证据判定是否属于请求范围；对不确定对象进入人工或治理复核，不自动删除。

“经历必留痕迹”也不是无条件守恒律。数字记录可被物理清除，工作态会自然释放，参数影响有时可通过重训或专门遗忘算法降低；另一方面，已被复制到无控制权系统中的内容可能无法完全追回。可达到何种程度取决于存储控制、复制历史、模型结构和验证能力，而不是由碳基或硅基材料共同遵循的几何法则决定。系统必须报告可控制域、不可控制域和剩余风险。

验收从范围证明开始。系统列出命中的原件、派生物、索引、缓存、备份与参数版本，执行声明操作后进行普通检索、对抗提示、成员推断、能力回归和恢复测试。查询未命中只证明特定接口下不可访问；文件不存在只证明该位置已清除；参数测试未恢复目标信息只能给出统计证据。最终删除证书应包含请求版本、授权、操作日志、测试覆盖、残余副本和无法保证的边界。由此，记忆消除的本体不再是“把容器形状抹平”，而是对有身份、有依赖、有权限且分布在多层载体中的对象执行可审计变更。
